% GENERATED FILE. Edit canonical lessons, then run npm run generate:books.
% canonical-source-hash: fnv1a64:8148e5f5d14f407c
% canonical-lessons: FR-C409-clignotant, FR-C409-phare, FR-C409-virage, FR-C409-chaussee, FR-C409-devier

\chapter{Dans les virages}
\label{ch:fr-dans-les-virages}
\hlchaptermodality{409}

\begin{chapteropening}
\textbf{By the end of this chapter:} \emph{I can talk about the road surface and signalling: the indicator, the headlights, the bend, the carriageway and turning aside.}

Complete the last lesson of chapter 409: I can talk about the road surface and signalling: the indicator, the headlights, the bend, the carriageway and turning aside.
\end{chapteropening}

\section[le clignotant]{le clignotant — the indicator, the turn signal}

\label{lesson:FR-C409-clignotant}



\begin{quote}
Before the new one: say the French for to fine, then the French for to inspect.
\end{quote}

\subsection*{You'll want to know: le clignotant}
\textbf{le clignotant} — \textquotedblleft{}the indicator, the turn signal\textquotedblright{}.

\begin{quote}
Il tourne sans mettre son clignotant.
\end{quote}

\begin{quote}
\emph{He turns without using his indicator.}
\end{quote}

Built on \emph{clignoter} 'to blink on and off'.

\begin{cousinweb}
\emph{clignotant} is built on \emph{clignoter} 'to flash on and off', from \emph{cligner} 'to blink'. The origin of \emph{cligner} is not settled, and EtymDB does not list the word.
\end{cousinweb}

\begin{grammarlens}[title={What you've built}]
One more word for telling, explaining and describing at B1.
\end{grammarlens}

\subsection*{Your turn}
\begin{itemize}
\raggedright
  \item \emph{Say it:} \emph{le clignotant}
  \item \emph{Say it:} \emph{le clignotant}, once more
  \item \emph{Say it:} a sentence of your own with \emph{le clignotant}
  \item \emph{Recall:} say the French for to fine, then the French for to inspect, then say \emph{le clignotant} again
\end{itemize}

\begin{tcolorbox}[breakable,colback=teal!4,colframe=teal!35!black,title={Before you move on}]
What does le clignotant mean? (\textquotedblleft{}The indicator, the turn signal\textquotedblright{}.) Say it once more.
\end{tcolorbox}
\section[le phare]{le phare — the headlight; the lighthouse}

\label{lesson:FR-C409-phare}



\begin{quote}
Before the new one: say the French for to inspect, then the French for the indicator.
\end{quote}

\subsection*{You'll want to know: le phare}
\textbf{le phare} — \textquotedblleft{}the headlight; the lighthouse\textquotedblright{}.

\begin{quote}
La nuit, les phares éclairent la route.
\end{quote}

\begin{quote}
\emph{At night the headlights light up the road.}
\end{quote}

\emph{Un phare} is also a lighthouse.

\begin{cousinweb}
\emph{phare} comes from Latin \emph{pharus} 'lighthouse', from Greek \emph{pháros}, first the name of an island whose lighthouse was one of the wonders of the ancient world (EtymDB).
\end{cousinweb}

\begin{grammarlens}[title={What you've built}]
One more word for telling, explaining and describing at B1.
\end{grammarlens}

\subsection*{Your turn}
\begin{itemize}
\raggedright
  \item \emph{Say it:} \emph{le phare}
  \item \emph{Say it:} \emph{le phare}, once more
  \item \emph{Say it:} a sentence of your own with \emph{le phare}
  \item \emph{Recall:} say the French for to inspect, then the French for the indicator, then say \emph{le phare} again
\end{itemize}

\begin{tcolorbox}[breakable,colback=teal!4,colframe=teal!35!black,title={Before you move on}]
What does le phare mean? (\textquotedblleft{}The headlight; the lighthouse\textquotedblright{}.) Say it once more.
\end{tcolorbox}
\section[le virage]{le virage — the bend, the curve (in a road)}

\label{lesson:FR-C409-virage}



\begin{quote}
Before the new one: say the French for the indicator, then the French for the headlight.
\end{quote}

\subsection*{You'll want to know: le virage}
\textbf{le virage} — \textquotedblleft{}the bend, the curve (in a road)\textquotedblright{}.

\begin{quote}
Le bus ralentit dans le virage.
\end{quote}

\begin{quote}
\emph{The bus slows down on the bend.}
\end{quote}

\emph{Un virage dangereux} — a dangerous bend.

\begin{cousinweb}
\emph{virage} is built on \emph{virer} 'to turn' with \emph{-age}. The origin of \emph{virer} is not settled; it is often linked to Latin \emph{gyrare} 'to turn round'.
\end{cousinweb}

\begin{grammarlens}[title={What you've built}]
One more word for telling, explaining and describing at B1.
\end{grammarlens}

\subsection*{Your turn}
\begin{itemize}
\raggedright
  \item \emph{Say it:} \emph{le virage}
  \item \emph{Say it:} \emph{le virage}, once more
  \item \emph{Say it:} a sentence of your own with \emph{le virage}
  \item \emph{Recall:} say the French for the indicator, then the French for the headlight, then say \emph{le virage} again
\end{itemize}

\begin{tcolorbox}[breakable,colback=teal!4,colframe=teal!35!black,title={Before you move on}]
What does le virage mean? (\textquotedblleft{}The bend, the curve (in a road)\textquotedblright{}.) Say it once more.
\end{tcolorbox}
\section[la chaussée]{la chaussée — the road surface, the carriageway}

\label{lesson:FR-C409-chaussee}



\begin{quote}
Before the new one: say the French for the headlight, then the French for the bend.
\end{quote}

\subsection*{You'll want to know: la chaussée}
\textbf{la chaussée} — \textquotedblleft{}the road surface, the carriageway\textquotedblright{}.

\begin{quote}
La chaussée est glissante après la pluie.
\end{quote}

\begin{quote}
\emph{The road surface is slippery after the rain.}
\end{quote}

The part of the street for cars, not the pavement.

\begin{cousinweb}
\emph{chaussée} is inherited from Late Latin \emph{calciāta} 'paved with lime', from \emph{calx} 'lime' (EtymDB). English \emph{causeway} is a cousin.
\end{cousinweb}

\begin{grammarlens}[title={What you've built}]
One more word for telling, explaining and describing at B1.
\end{grammarlens}

\subsection*{Your turn}
\begin{itemize}
\raggedright
  \item \emph{Say it:} \emph{la chaussée}
  \item \emph{Say it:} \emph{la chaussée}, once more
  \item \emph{Say it:} a sentence of your own with \emph{la chaussée}
  \item \emph{Recall:} say the French for the headlight, then the French for the bend, then say \emph{la chaussée} again
\end{itemize}

\begin{tcolorbox}[breakable,colback=teal!4,colframe=teal!35!black,title={Before you move on}]
What does la chaussée mean? (\textquotedblleft{}The road surface, the carriageway\textquotedblright{}.) Say it once more.
\end{tcolorbox}
\section[dévier]{dévier — to turn aside, to divert}

\label{lesson:FR-C409-devier}



\begin{quote}
Before the new one: say the French for the bend, then the French for the road surface.
\end{quote}

\subsection*{You'll want to know: dévier}
\textbf{dévier} — \textquotedblleft{}to turn aside, to divert\textquotedblright{}.

\begin{quote}
La voiture a dévié de sa route sur le pont.
\end{quote}

\begin{quote}
\emph{The car swerved off course on the bridge.}
\end{quote}

Also \emph{dévier la circulation} — to divert the traffic.

\begin{cousinweb}
\emph{dévier} is a learned borrowing from Late Latin \emph{dēviāre} 'to leave the road', from \emph{via} 'way' (EtymDB).
\end{cousinweb}

\begin{grammarlens}[title={What you've built}]
One more word for telling, explaining and describing at B1.
\end{grammarlens}

\subsection*{Your turn}
\begin{itemize}
\raggedright
  \item \emph{Say it:} \emph{dévier}
  \item \emph{Say it:} \emph{dévier}, once more
  \item \emph{Say it:} a sentence of your own with \emph{dévier}
  \item \emph{Recall:} say the French for the bend, then the French for the road surface, then say \emph{dévier} again
\end{itemize}

\begin{tcolorbox}[breakable,colback=teal!4,colframe=teal!35!black,title={Before you move on}]
What does dévier mean? (\textquotedblleft{}To turn aside, to divert\textquotedblright{}.) Say it once more.
\end{tcolorbox}
